\documentclass{resume}
\usepackage{zh_CN-Adobefonts_external}
\usepackage{linespacing_fix}
\usepackage{cite}
\usepackage{hyperref}
\hypersetup{
colorlinks=true,
linkcolor=cyan,
filecolor=magenta,
urlcolor=blue,
}

\begin{document}
\pagenumbering{gobble}

%***"%"后面的所有内容是注释而非代码，不会输出到最后的PDF中
%***使用本模板，只需要参照输出的PDF，在本文档的相应位置做简单替换即可
%***修改之后，输出更新后的PDF，只需要点击Overleaf中的"Recompile"按钮即可

%在大括号内填写其他信息，最多填写4个，但是如果选择不填信息，
%那么大括号必须空着不写，而不能删除大括号。
%\otherInfo后面的四个大括号里的所有信息都会在一行输出
%如果想要写两行，那就用两次这个指令（\otherInfo{}{}{}{}）即可

%***********个人信息**************
\MyName{段圣宣}
\sepspace
\SimpleEntry{齐齐哈尔大学}
\SimpleEntry{2598438968@qq.com}
\SimpleEntry{17771863696}

%***********教育背景**************
\section{教育背景}

\datedsubsection{\textbf{齐齐哈尔大学}，数据科学与大数据技术，\textit{本科}}{2023.9 - 2027.6}

%***********实习经历**************
\section{实习经历}

\datedsubsectionlarge{\textbf{上海稳博投资｜量化开发工程师实习生}}{2026.1 - 2026.3}
\begin{itemize}[parsep=0.3ex]
\item \large 滚木
\end{itemize}

\datedsubsectionlarge{\textbf{波克私募｜量化开发工程师实习生}}{2026.4 - 至今}
\begin{itemize}[parsep=0.3ex]
\textbf{流式因子计算引擎}
\item \large 设计 20+ 个高性能滚动窗口算子，使用 $\sqrt{n}$ 分块有序结构替代堆分配 BST，并在小窗口场景采用 AVX-512 暴力扫描，利用连续内存、低分支和硬件预取优化热路径。
\item \large 优化因子DAG计算图并行调度算法，根据数据特性优化矩阵算子执行路径，根据缓存时间和空间局部性优化计算执行路径，延迟降低高达98\%。
\item \large 实现因子计算引擎中代数化简与图优化模块。在 DAG 计算图上实现公共子表达式消除、常量折叠、公共分母提取等优化，图计算节点数下降30\%，
在 SIMD 流水线上充分并行；整体计算吞吐较未化简基线提升约 53\%。
\textbf{低延迟量化交易执行与策略事件引擎}
\item \large 基于 C++20 开发的量化交易策略事件引擎，通过共享内存连接实时因子计算系统与订单管理系统，实现多频数据对齐、策略事件驱动、历史数据查询及订单状态同步，并提供 Python 接口，支持策略快速开发与运行。
\item \large 设计支持并发查询修改订单状态的数据结构和数据同步机制，使用二级索引结构优化订单查询；使用追加写入的块状链表保存记录编号，并通过提交位置控制读写可见性，保证缓存良好性和低延迟，写入p50延迟低至3us，读取低至1us。

\end{itemize}


%***********奖励荣誉**************
\section{奖励荣誉}

{\footnotesize
\datedsubsection{ACM-ICPC 国际大学生程序设计竞赛亚洲区域赛(武汉)银奖}{2025.11}

\datedsubsection{CCPC 中国大学生程序设计竞赛全国邀请赛(南昌)冠军}{2025.09}

}

%***********专业技能**************
\section{专业技能}

{\footnotesize
\datedsubsection{拥有较强的算法设计和实现能力，掌握各类数据结构和算法优化方案}{}

\datedsubsection{熟悉现代 C++（C++17/20），熟悉 RAII、移动语义，理解虚函数表布局、并发编程、内存模型与无锁数据结，熟悉 STL 容器底层实现}{}

\datedsubsection{熟悉 CPU 体系结构与性能优化：理解流水线、乱序执行、分支预测、NUMA 与 cache 多级结构；熟悉 CPU 缓存层次结构，能针对性设计缓存友好的数据布局以减少访存开销；熟练使用 perf、火焰图等工具进行性能与内存剖析。}{}

\datedsubsection{熟悉操作系统基础原理，理解进程与线程、虚拟内存、页式管理、锁与同步、系统调用等底层机制。}{}

}

\end{document}
